% ===== PRÉAMBULE CANONIQUE — à coller VERBATIM en tête de CHAQUE .tex =====
\documentclass[11pt,a4paper]{article}
\usepackage[utf8]{inputenc}\usepackage[T1]{fontenc}\usepackage{lmodern}
\usepackage[french]{babel}
\usepackage{amsmath,amssymb,mathtools}\usepackage{icomma}
\usepackage[version=4]{mhchem}          % \ce{...} : équations & formules
\usepackage{siunitx}\sisetup{locale=FR} % \qty{}{}, \si{}, \ang{}
\usepackage{chemfig}                    % \chemfig{...} : structures développées
\usepackage{xcolor}\usepackage{colortbl}
\usepackage{tikz}\usepackage{pgfplots}\pgfplotsset{compat=1.18}
\usepackage[most]{tcolorbox}\usepackage{enumitem}\usepackage{array,booktabs}
\usepackage[a4paper,margin=2.2cm]{geometry}
\usetikzlibrary{arrows.meta,calc,angles,quotes,positioning,patterns,backgrounds,shapes.geometric}
\definecolor{primary}{RGB}{222,49,99}\definecolor{primarydark}{RGB}{150,22,55}
\definecolor{defcol}{RGB}{0,114,178}\definecolor{methcol}{RGB}{52,92,148}
\definecolor{excol}{RGB}{0,135,90}\definecolor{attncol}{RGB}{200,40,55}
\tcbset{boxbase/.style={enhanced,breakable,boxrule=0.9pt,arc=2.5pt,
  left=3mm,right=3mm,top=2.3mm,bottom=2.3mm,fonttitle=\bfseries,coltitle=white}}
% --- boîtes TUTORIEL (noms EXACTS, ne pas renommer) ---
\newtcolorbox{cle}[1][]{boxbase,colback=defcol!6,colframe=defcol,
  title={\textsc{À retenir}\ifx\relax#1\relax\relax\else\textmd{ : \itshape #1}\fi}}
\newtcolorbox{methode}[1][]{boxbase,colback=methcol!7,colframe=methcol,sharp corners,
  title={\textsc{Méthode}\ifx\relax#1\relax\relax\else\textmd{ --- \itshape #1}\fi}}
\newtcolorbox{exemplebox}[1][]{boxbase,colback=excol!5,colframe=excol,
  title={\textsc{Exemple}\ifx\relax#1\relax\relax\else\textmd{ : \itshape #1}\fi}}
\newtcolorbox{piege}[1][]{boxbase,colback=attncol!6,colframe=attncol,
  title={\textsc{Piège}\ifx\relax#1\relax\relax\else\textmd{ : \itshape #1}\fi}}
\newtcolorbox{remarque}[1][]{boxbase,colback=black!4,colframe=black!45,
  title={\textsc{Remarque}\ifx\relax#1\relax\relax\else\textmd{ : \itshape #1}\fi}}
% --- boîtes EXERCICE / CORRIGÉ (noms EXACTS, auto-comptées) ---
\newtcolorbox[auto counter]{exo}[1][]{boxbase,colback=primary!4,colframe=primary,
  title={\textsc{Exercice}~\thetcbcounter\ifx\relax#1\relax\relax\else\textmd{ --- \itshape #1}\fi}}
\newtcolorbox[auto counter]{corrige}[1][]{boxbase,colback=excol!4,colframe=excol,
  title={\textsc{Corrigé}~\thetcbcounter\ifx\relax#1\relax\relax\else\textmd{ --- \itshape #1}\fi}}
\newcommand{\R}{\mathbb{R}}\newcommand{\N}{\mathbb{N}}\newcommand{\Z}{\mathbb{Z}}\newcommand{\ds}{\displaystyle}
% (un \usepackage{fancyhdr}+en-tête est possible ; il est ignoré côté web)
% =========================================================================
\begin{document}
\begin{corrige}[octet et duet]
\begin{enumerate}[label=\alph*)]
  \item $8$ électrons de valence (un octet : couche de valence saturée).
  \item L'hélium en a $2$ : c'est un \textbf{duet} (sa couche K est saturée).
  \item Leur couche de valence est saturée : ils n'ont ni tendance à gagner ni à perdre d'électrons,
        donc ils sont stables et quasi inertes.
\end{enumerate}
\end{corrige}

\begin{corrige}[reconnaître une famille]
\begin{enumerate}[label=\alph*)]
  \item $1$ e$^-$ de valence : \textbf{alcalins} (forment $\ce{X^{+}}$).
  \item $2$ e$^-$ de valence : \textbf{alcalino-terreux} (forment $\ce{X^{2+}}$).
  \item $7$ e$^-$ de valence : \textbf{halogènes} (forment $\ce{X^{-}}$).
  \item $8$ e$^-$ de valence : \textbf{gaz nobles} (saturés, inertes).
\end{enumerate}
\end{corrige}

\begin{corrige}[quel ion se forme ?]
Les métaux perdent leurs électrons de valence ; les non-métaux en gagnent pour compléter l'octet.
\begin{enumerate}[label=\alph*)]
  \item $\ce{Na}$ perd $1$ e$^-$ $\to \ce{Na^{+}}$.
  \item $\ce{Mg}$ perd $2$ e$^-$ $\to \ce{Mg^{2+}}$.
  \item $\ce{O}$ gagne $8-6 = 2$ e$^-$ $\to \ce{O^{2-}}$.
  \item $\ce{Cl}$ gagne $8-7 = 1$ e$^-$ $\to \ce{Cl^{-}}$.
\end{enumerate}
\end{corrige}

\begin{corrige}[sens des tendances périodiques]
\begin{enumerate}[label=\alph*)]
  \item L'électronégativité augmente vers la \textbf{droite} et vers le \textbf{haut}.
  \item Le rayon atomique augmente vers la \textbf{gauche} et vers le \textbf{bas} (sens inverse).
  \item L'énergie d'ionisation varie dans le même sens que \textbf{l'électronégativité}
        (vers la droite et le haut).
\end{enumerate}
\end{corrige}

\begin{corrige}[comparer deux éléments]
\begin{enumerate}[label=\alph*)]
  \item \textbf{$\ce{Cl}$} : plus à droite dans la même période, donc plus électronégatif que $\ce{Na}$.
  \item \textbf{$\ce{K}$} : plus bas dans la même colonne, donc plus grand rayon que $\ce{Li}$.
  \item \textbf{$\ce{F}$} : plus à droite dans la même période, donc énergie d'ionisation plus grande
        que $\ce{Li}$.
\end{enumerate}
\end{corrige}

\end{document}
